\section{Methods}
\subsection{Leaky integrate-and-fire neuron model}
We employed the leaky integrate-and-fire (LIF) spiking neuron model to describe the membrane potential dynamics of neurons:
\begin{equation}
    \dfrac{dV_i}{dt}= -LV_i(t) + I_i(t),
    \label{eq:LIF}
    \end{equation} 
where the sub-threshold membrane potential $V_i(t)$ of a neuron $i$ is driven by the total synaptic current $I_i(t)$ and $L=0.05$ \si{\per\milli\second} is the leak conductance. 
When the membrane potential $V_i(t)$ exceeds a threshold $V_{\rm th}=20$ \si{\milli\volt} a spike is emitted, as represented by a Dirac delta function. Afterward, the membrane potential $V_i(t)$ is reset to the resting potential $V_{\rm res}=0$ mV, followed by a refractory period $T_{\rm ref}=5$ ms. The synaptic current takes the form
\begin{equation}
    I_i(t)= \sum_{j}w_{ij}S_j(t)+I_i^{\rm ext}(t),
    \label{eq:current}
\end{equation}
where $S_j(t)=\sum_k \delta(t-t^k_j)$ represents the spike train generated by presynaptic neurons.

A final output $\mathbf{y}$ is readout from the spike count $\mathbf{n}(\Delta t)$ of a population of spiking neurons over a time window of duration $\Delta t$ as follows
\begin{equation}
y_i(\Delta t) = 
\frac{1}{\Delta t}\sum_{j} w_{ij}n_{j}(\Delta t) + \beta_i, %NB decoder has a bias
\label{eq:snn_readout}
\end{equation}
where $w_{ij}$ and $\beta_i$ are the weights and biases of the readout, respectively. 
A key characteristic of the readout is that its variance decreases as the time window $\Delta t$ increases.

\subsection{Moment embedding for the leaky integrate-and-fire neuron model}
The moment embedding approach~\cite{Feng2006, Lu2010, Qi2023} begins with mapping the fluctuating activity of spiking neurons to their respective first- and second-order moments
\begin{equation}
\mu_i = \lim_{\Delta t\to\infty} \dfrac{\mathbb{E}[n_i(\Delta t)]}{\Delta t}, 
\label{eq:def_mean}
\end{equation}
and
\begin{equation}
\Sigma_{ij} = \lim_{\Delta t\to\infty} \dfrac{{\rm Cov}[n_i(\Delta t),n_j(\Delta t)]}{\Delta t}, 
\label{eq:def_cov}
\end{equation}
where $n_i(\Delta t)$ is the spike count of neuron $i$ over a time window $\Delta t$. 
In practice, the limit of $\Delta t\to\infty$ is interpreted as a sufficiently large time window relative to the timescale of neural fluctuations. 
We refer to the moments $\mu_i$ and $\Sigma_{ij}$ as the mean firing rate and the firing covariability in the context of MNN, respectively. 

For the LIF neuron model [equation (\ref{eq:LIF})], the statistical moments of the synaptic current are equal to~\cite{Feng2006, Lu2010}
\begin{numcases}{}
\hat{\mu}_i = \textstyle\sum_kw_{ik}\mu_k + \hat{\mu}_i^{\rm ext},\label{eq:sum_mean}
\\
\hat{\Sigma}_{ij}=\textstyle\sum_{kl} w_{ik}C_{kl}w_{jl} + \hat{C}_{ij}^{\rm ext},\label{eq:sum_cov}
\end{numcases}
where $w_{ik}$ is the synaptic weight and $\hat{\mu}_i^{\rm ext}$ and $\hat{\Sigma}_{ij}^{\rm ext}$ are the mean and covariance of an external current, respectively.  
Note that from equation (\ref{eq:sum_cov}), it becomes evident that the synaptic current is correlated even if the presynaptic spike trains are not. 
Next, the first- and second-order moments of the synaptic current are mapped to that of the spiking activity of the post-synaptic neurons. For the LIF neuron model, this mapping can be obtained in closed form through a mathematical technique known as the diffusion approximation~\cite{Feng2006, Lu2010} as
\begin{numcases}{}
\mu_i = \phi_\mu(\bar{\mu}_i,\bar{\sigma}_i),\label{eq:ma_mu}\\
\sigma_i = \phi_\sigma(\bar{\mu}_i,\bar{\sigma}_i),\label{eq:ma_sig}\\
\rho_{ij} = \chi(\bar{\mu}_i,\bar{\sigma}_i)\chi(\bar{\mu}_j,\bar{\sigma}_j)\bar{\rho}_{ij},\label{eq:ma_chi}
\end{numcases}
where the correlation coefficient $\rho_{ij}$ is related to the covariance as $\Sigma_{ij}=\sigma_i\sigma_j\rho_{ij}$. The mapping given by equation (\ref{eq:ma_mu})-(\ref{eq:ma_chi}) is called the moment activation, which is differentiable so that gradient-based learning algorithms can be implemented and the learning framework is known as the moment neural network (MNN)~\cite{Qi2023}.



\subsection{The motion direction detection task}
\subsubsection{Model setup for training}
We trained the MNN model for motion direction detection as detailed in our previous work~\cite{Qi2023}.
The network configuration includes 1054 sensory neurons, 1054 hidden neurons, and 2 readout neurons. 
In the input layer, 527 pairs of intensity and change detectors are evenly distributed across a unit hexagonal grid. 
The sensory neurons, characterized by input response coefficients \(\alpha = \beta = 1\), the temporal angular frequency of \(\omega = 1\), the observation window  \(\Delta t = 1\) and the spatial frequency length \(k = |\mathbf{k}| = 5\pi\), follow our proposed encoding scheme  [equation~(\ref{eq:vo_input_mean})-(\ref{eq:vo_input_cov})] for generating inputs.
The hidden layer comprises a synaptic summation layer, succeeded by a moment batch normalization layer and a moment activation layer. 
After training, the parameters of the moment batch normalization layer were integrated into the synaptic summation layer for optimization.

For effective training of the model parameters \(\bm{\Phi}\), we introduced a loss function based on the readout mean \(\hat{\mu}\) and covariance \(\hat{\Sigma}\),
\begin{equation}
    L(\bm{\Phi}) = \int p(\mathbf{x|\bm{\Phi}}) \arccos\left( \frac{\mathbf{x} \cdot \mathbf{t}}{|\mathbf{x}||\mathbf{t}|}\right) d\mathbf{x} = \mathbb{E}\left[\arccos\left( \frac{\mathbf{x} \cdot \mathbf{t}}{|\mathbf{x}||\mathbf{t}|}\right) \right],
    \label{eq:loss_function}
\end{equation}
where \(\mathbf{x} \sim \mathbb{N}(\hat{\mu}, \hat{\Sigma}|\bm{\Phi})\) represents a Gaussian-distributed readout and \(\mathbf{t}\) is the coordinates of ground truth (\(\cos \theta, \sin \theta\)).
Due to the complexity of this loss function, we approximated it using a finite-sample method. 
Specifically, we estimated \(L(\bm{\Phi})\) as \(\frac{1}{N}\sum_{n=1}^N \arccos (\mathbf{x}^n\mathbf{t}/|\mathbf{x}^n||\mathbf{t}|)\), where each \(\mathbf{x}^n\) is generated using Cholesky decomposition \(\Sigma =  LL^T\), express as \(\mathbf{x}^n = L\mathbf{z}^n + \mathbf{\mu}\) with \(\mathbf{z}^n\) being an uncorrelated unit normal random variable.

The model was trained for 150 epochs using the AdamW optimizer with its default hyperparameters (0.001 learning rate and 0.01 weight decay). 
During each epoch, 10,000 samples were created by randomly selecting contrasts (from 0 to 0.8) and motion directions (from \(-\pi\) to \(\pi\)) to guarantee a comprehensive and varied training dataset.

\subsubsection{Model setup for SNN simulation}
The trained MNN parameters were directly utilized to reconstruct the SNN, as detailed in our previous work~\cite{Qi2023}. For the motion direction detection task, we maintained a consistent contrast level \(c = 0.8\) and selected 50 motion directions uniformly distributed between \(-\pi\) and \(\pi\).
The membrane potentials of hidden neurons were initially set to random values between the resting potential\(V_{\rm res}\) and the threshold  \(V_{\rm th}\). 
To ensure accurate calculation of mean spike rates and trial-to-trial covariance for decoding, we conducted 500 trials for each direction, simulating each for 1256 ms at a time increment of \(\delta t = 0.1\) ms. 
The task inputs were generated via the inhomogeneous Poisson process as described in section \ref{sec:vo_encode}.

We collected responses from the hidden neurons and readouts from the networks at each time step for analysis. 
It is crucial to differentiate these SNN readouts from those obtained from the MNNs. 
The SNN readouts were accumulated over time, and we calculated the discrepancy between the estimated direction represented by these readouts and the actual motion directions.
\subsubsection{Population activity analysis}
We began by tallying the spikes of all sensory and hidden neurons within 1 ms intervals throughout the duration of the simulation. 
This produced a vector representing the population firing rate at each time point. 
The relative deviation of the population spike count was calculated as:
\begin{equation}
D(t) = \frac{R(t) - \langle R \rangle}{\langle R \rangle}, 
\end{equation}
where \(R(t)\) is the population spike count and \(\langle R \rangle\) is its trial-averaged mean rate. 
Using the relative deviation \(D(t)\), we computed the power spectrum density and autocorrelation (as shown in Fig.~\ref{fig:visual_orientation_snn}\textbf{c}), averaging their mean and standard deviation across 50 directions, each estimated with 500 trials.

For further analysis, we used the firing rates of sensory neurons and membrane potentials of hidden neurons recorded from 100 ms to 220 ms in the simulations. 
Data preprocessing involved subtracting the mean value from each trial. 
The Hilbert transform was then applied to calculate the instantaneous phase (\(\psi_i(t)\)) of each neuron. 
The average phase \(\bar \psi(t)\) and the instantaneous Kuramoto order parameter \(r(t)\) determined as follows:
\begin{equation}
r(t) e^{i\bar\psi(t)} = \frac{1}{N}\sum^{N}_{j=1} e^{i \psi_j(t)}. 
\end{equation}
The mean of \(r(t)\) was taken as a measure of synchrony during the simulation.
Since sensory neurons exhibited consistent firing rates in all trials, only the average direction was used for visualization. 
For hidden neurons, the mean and standard deviation of the average phase and the Kuramoto order parameter were first estimated through 500 trials for each direction and then averaged over all 50 directions.
\subsubsection{Information-theoretic analysis}
The mutual information between the motion direction \(\theta\) and the readout \(\mathbf{x}\) is defined in terms of the entropy of all possible responses and conditional responses as
\begin{equation}
    I_{\rm tot}(\mathbf{x}; \theta) = h(\mathbf{x}) - h(\mathbf{x}|\theta).
\end{equation}
The entropy \(h(\mathbf{x})\) and the conditional entropy \(h(\mathbf{x}| \theta)\) are given by
\begin{align}
    h(\mathbf{x}) &= - \int p(\mathbf{x}) \log p(\mathbf{x}) d\mathbf{x},\\
    h(\mathbf{x}|\theta) &= -\sum_i p(\theta_i) \int p(\mathbf{x}|\theta_i) \log p(\mathbf{x}|\theta_i) d\mathbf{x},
\end{align}
where \( p(\mathbf{x}|\theta)\) is assumed to be a Gaussian distribution with readout mean \(\hat{\mathbf{\mu}}\) and covariance \(\hat{\Sigma}\) condition on the stimulus \(\theta\).
The mean \(\hat{\mu}\) covariance \(\hat{\Sigma}\) of the readout \(\mathbf{x}\) within readout time \(\Delta t\) were estimated through accumulative readouts, using data from SNN simulations.  
Assuming a uniform prior and discretizing motion directions into 50 bins (matching the number of directions used for SNN simulations), the distribution of overall neural responses was calculated as follows:
\begin{equation}
    p(\mathbf{x}) = \sum_i p(\theta_i)p(\mathbf{x}|\theta_i).
\end{equation}
We used the information breakdown analysis~\cite{Pola2003, montani2007role} to further dissect the mutual information \(I_{\rm tot}\) into three components, allowing us to assess the contributions of individual neuron responses \(I_{\rm lin}\), signal similarity among neurons \(I_{\rm sigsim}\), and the correlation in neuron activities \(I_{\rm cor}\).
\(I_{\rm lin}\) measures the total amount of information that would be transmitted if all neurons were independent, which is given by:
\begin{equation}
    I_{\rm lin} = \sum_j [h(x_j) - h(x_j | \theta)],
\end{equation}
where \(x_j\) is the possible readout of neuron \(j\).
\(I_{\rm sigsim}\) quantifies the information loss arising from redundancy due to overlap in the tuning curves about response \(x_j\), which is given by:
\begin{equation}
    I_{\rm sigsim} = h(\mathbf{x}_{\rm ind}) - \sum_j h(x_j),
\end{equation}
where the independent population response \(\mathbf{x}_{\rm ind}\) is defined as:
\begin{equation}
    p(\mathbf{x}_{\rm ind} | \theta) = \prod_i p(x_i | \theta).
\end{equation}
The last component \(I_{\rm cor}\) accounts for the rest part of \(I_{\rm tot}\), quantifies the total amount of information due to the correlated activity in the neural responses:
\begin{equation}
    I_{\rm cor} = I_{\rm tot} - I_{\rm lin} - I_{\rm sigsim}
\end{equation} 
Since there is no analytical expression for the entropy of unconditional neural response \(p(x)\), we estimated this information through Monte Carlo sampling.
The mean \(\hat{\mathbf{\mu}}\) and covariance \(\hat{\Sigma}\) were computed based on accumulative readouts for all directions in a readout time window \(\Delta t\) (Fig~\ref{fig:visual_orientation_snn}\textbf{e}). 
Subsequently, we sampled 10000 data points based on these \(\hat{\mu}\) and \(\hat{\Sigma}\) at time \(t\) to empirically estimate the information as a function of the readout time window \(\Delta t\).
The theoretic bound of \(I_{\rm tot}\) is the entropy of direction, where \(H(\theta) = \sum_i p(\theta_i)\log \theta_i = \log 50\) nats and we defined the time when the readout information converges as the first time when \(|I_{\rm tot} - H(\theta)| < 0.1\) nats.

\subsection{The fine-grain classification task}
\subsubsection{Model setup for training}
For our fine-grain classification task, we utilized the Caltech-UCSD Birds-200-2011 dataset \cite{WahCUB_200_2011}, comprising 11,788 images of 200 bird species. 
Adhering to the recommended dataset split, we divided it into a training set of 5,994 images and a test set of 5,794 images.

Our image preprocessing was in agreement with previous work~\cite{Wang2020}. 
We resized the images to have a shorter side of 448 pixels, then center-cropped them to \(448 \times 448\) resolution. 
To augment the training data, the images were randomly flipped horizontally. Normalization was applied using z-score normalization, with RGB channel means (0.485, 0.456, 0.406) and standard deviations (0.229, 0.224, 0.225) derived from the ImageNet dataset~\cite{simonyan2014very}.

Feature maps were extracted using a pretrained VGG16 model~\cite{simonyan2014very}  from PyTorch, with its parameters fixed. The resulting feature maps, of dimensionality \(512 \times 28 \times 28\), were flattened and processed to calculate channel-wise mean and covariance [equation~(\ref{eq:theory_define_mean})-(\ref{eq:theory_define_cov})], serving as inputs for our MNN model.

To evaluate our theory, we modified the inputs in various ways. 
The \textit{With corr} model received both mean and covariance as inputs, while the \textit{Without corr} model had its covariance off-diagonal entries zeroed, erasing correlation information. 
The \textit{ANN} was trained using only the mean. Each model, including the MNNs and ANN, consisted of an input layer (512 neurons), two hidden layers (1024 neurons each), and a readout layer (200 neurons). 
They were trained 5 times, each time for 150 epochs, using standard cross-entropy loss and AdamW optimizer with the default setting.

\subsubsection{Model setup for SNN simulation}
We reconstructed the SNNs using the parameters of the trained MNNs under both \textit{With corr} and \textit{Without corr} conditions. 
To evaluate the performance of the SNNs on the fine-grain classification task, we simulated the network for each image in the test set, each for 100 trials, with each simulation running for 150 ms at a time increment of \(\delta t = 0.1\) ms. 
The task inputs were generated by sampling from a Gaussian distribution, using the mean and covariance of the respective input conditions. 

During these simulations, we accumulated the readouts and determined the prediction of the model by identifying the dimension index with the highest cumulative readouts. 
The accuracy of the prediction was determined by taking the average of the proportion of correct predictions made for all the samples in the test set in all the trials.